\hypertarget{DURUM_OZETI}{%
\chapter{Durum özeti --- Faz 3, parça 2 sonu (18 Ekim 2026)}\label{DURUM_OZETI}}

Yeni bir sohbette devam etmek için ana prompt ile bu dosyayı birlikte verin. Ayrıntılar: \belgebag{nace_kartlar_13-15}{}{\kod{nace\_\allowbreak{}kartlar\_\allowbreak{}13-\allowbreak{}15.\allowbreak{}md}}, \belgebag{nace_kartlar_10-12}{}{\kod{nace\_\allowbreak{}kartlar\_\allowbreak{}10-\allowbreak{}12.\allowbreak{}md}}, \belgebag{02c_hazirlik}{}{\kod{02c\_\allowbreak{}hazirlik.\allowbreak{}md}}, \belgebag{02b_secenekler_hukuk}{}{\kod{02b\_\allowbreak{}secenekler\_\allowbreak{}hukuk.\allowbreak{}md}}, \belgebag{02a_politika_uzayi}{}{\kod{02a\_\allowbreak{}politika\_\allowbreak{}uzayi.\allowbreak{}md}}, \belgebag{01_teshis}{}{\kod{01\_\allowbreak{}teshis.\allowbreak{}md}}.

\textbf{Çerçeve (değişmedi).} Çıkış yönü veri; AB hasmane; üyelik perspektifi yok. İki takvim: T2028 ve T2029. Önerilen bitiş durumu S2.

\textbf{Faz 2 (özet).} Yerliye özgü ÖTV indirimi GATT III:2 ihlali; eşdeğeri bütün menşelere eşit ÖTV/KDV takası. Genel dengede S4'te AB'ye ihracat −\%19, S2'de −\%11. 51 sessiz hamle; son tarihler XXVIII:5 31.12.2026, DS629 24.11.2026.

\textbf{Faz 3 altyapısı (230 sınıf).} Ticaret, yapı, PCI, AB tarifesi, Monte Carlo senaryoları, skorlar, ana matris. Yeni: PEM kümülasyonu kanalı (yalnız S4), bölge tablosu, parçalar arası kümülatif defter. İmalat geneli (geçici): S4'te istihdam −46 bin (P10 −70 bin), TÜFE 69 bp; S2'de −24 bin, 11 bp.

\textbf{Parça 1: 10--12 (33 sınıf).} Asıl risk 1/98'in askıya alınması (olasılık 0,25): AB'ye ihracat −\$1,7 mlr. S4 medyan etkileri küçük.

\textbf{Parça 2: 13--15 (21 sınıf).}

\begin{itemize}
\tightlist
\item
  1,24 milyon çalışan (imalatın \%25'i); ihracat \$35,9 mlr, \%48'i AB'ye; tamamı GB sanayi rejiminde. S4'te AB tarifesi \%10,2, giyimde \%11,8.
\item
  S4: istihdam −39 bin (P10 −61 bin), imalattaki S4 kaybının \%85'i; dış giyim −11,9 bin, iç giyim −9,4 bin. Cari denge −\$2,6 mlr. S2'de −14 bin.
\item
  Yeni kanal: PEM ortaklarına (Fas, Tunus, Balkanlar) kumaş ihracatı \$1,7 mlr, AB'ye dönen maruziyet \$1,1 mlr. S4'te kümülasyon düşüyor (PEM md. 8): −\$0,67 mlr, −4,8 bin istihdam. E-13-01 ilk müzakere talebi.
\item
  S2'de menşe: çift dönüşüm; Asya kumaşlı giyim (\%10--15) tercih dışı. Geri ödeme yasağı yeni maliyet değil (TEV bugün de var). Hücum: menşe kotası, örme giyimde baskı seçeneği.
\item
  TÜFE: c2 ile S4 şoku 15,6 bp; giyim ve ayakkabıda KDV takasıyla 3,6 bp. Takasın net maliyeti 17,5 mlr TL/yıl.
\item
  Bölgeler: Mardin--Batman--Şırnak--Siirt'te 13--15 imalat istihdamının \%78'i.
\item
  Duruşlar: 11 Geliştir, 2 Dönüştür (dış giyim, tabaklama), 3 Yönetilen uyum (iç giyim, diğer giyim, triko), 5 Koru.
\item
  Olumsuz bulgular: S4'te giyim kaybı önlenemiyor, ancak yumuşatılabiliyor. Kümülasyon tahmini üst sınır (κ, λ \etiketH{}).
\end{itemize}

\textbf{Kümülatif defter (10--15, 54 sınıf).} TÜFE S2 2,4, S4 8,2 bp (50 bp bütçe). Mali S2 1,65, S4 22,8 mlr TL/yıl (\%0,03 GSYH). Cari denge S4 P50 −\$2,75 mlr. İstihdam S4 P50 −40 bin.

\textbf{Düzeltmeler.}

\begin{itemize}
\tightlist
\item
  D-4: Tüzük 2026/2102 sığır derisini EUDR'den çıkardı; parça 1 metni düzeltildi.
\item
  D-5: parça 1'in S2 sayıları çok az değişti.
\item
  D-6: Faz 2'nin giyim tahmini (−59 bin) geçersiz.
\end{itemize}

\textbf{Açık konular.}

\begin{itemize}
\tightlist
\item
  Veri boşlukları: VB-02, 04, 06, 08--12, 14--17, 21--27, 30--39. Önemlileri: VB-36 firma kumaş menşei; VB-37 kümülasyon maruziyeti.
\item
  VB-32 (bağlı oranlar) 13--15'te bağlayıcı değil.
\item
  23.11 ve 23.14'te istihdam payı tanımsız; parça 5'te bakılacak.
\end{itemize}

\textbf{Sonraki adım.} Faz 3 parça 3: 16 Ağaç, 17 Kağıt, 18 Basım (EUDR ahşap ve kâğıt; SKDM yok).

\begin{gbilerleme}

İLERLEME: Faz 3 parça 2 / 16 tamamlandı (54 / 230 sınıf) --- SONRAKİ ADIM: Faz 3 parça 3 (16--18); DEVAM bekleniyor.

\end{gbilerleme}
